\documentclass[onecolumn,showpacs,preprintnumbers,amsmath,amssymb,superscriptaddress,preprint]{revtex4-1}
\newenvironment{cfigure}[1][tbp]{\begin{figure}[#1]\centering}{\end{figure}}
\newenvironment{cfigure1c}[1][tbp]{\begin{figure*}[#1]\centering}{\end{figure*}}

\def\cs{$^{137}$Cs~}
\def\ces{$^{134}$Cs~}
\def\rb{$^{87}$Rb~}
\def\cof{$^{57}$Co~}
\def\fe{$^{55}$Fe~}
\def\mn{$^{54}$Mn~}
\def\cos{$^{60}$Co~}
\def\halft{$t_{1/2}$~}
\def\cpd{$counts/(keV \cdot kg \cdot day)$~}
\def\naitl{NaI(T$\ell$)~}
\def\csitl{CsI(T$\ell$)~}
\def\PREP{Phys. Rep.}

\def\Journal#1#2#3#4{{#1} {\bf #2} (#3) #4}

% Some useful journal names
\def\NCA{Nuovo Cimento}
\def\NIM{Nucl. Instrum. Methods}
\def\NIMA{Nucl. Instrum. Methods A}
\def\NPB{Nucl. Phys. B}
\def\PLB{Phys. Lett.  B}
\def\PRL{Phys. Rev. Lett.}
\def\PRD{Phys. Rev. D}
\def\ZPC{Z. Phys. C}
\def\ASP{Astropart. Phys.}
\def\JKPS{J. Kor. Phys. Soc.}
\def\etal{{\it et al.}}


%\usepackage{elsart}
\usepackage{epsfig}
\usepackage{graphicx}
\usepackage{watermark}
\usepackage{comment}
\usepackage{multirow}
\usepackage{dcolumn}% Align table columns on decimal point
\usepackage{bm}% bold math
%\usepackage{lineno}
\usepackage{setspace}
\usepackage{graphicx}% Include figure files
\includecomment{printrobustnotes}
\includecomment{printallnotes}
\usepackage{xspace}
\usepackage{hyperref}
\hypersetup{
    colorlinks=true,
    linkcolor=blue,
    filecolor=magenta,      
    urlcolor=cyan,
}

\usepackage{listings} % Code listing
\usepackage{kotex} % Support Korean
\setmainhangulfont{NanumMyeongjo}

\usepackage{dcolumn}% Align table columns on decimal point
\usepackage{bm}% bold math 
\begin{document}

\begin{flushright}
CUP/INFRA/IT/1\\ %CUP/Experiments/Analysis/Note Number
Version 1.0\\ % Note version
\today
\end{flushright}

\begin{center}
  \includegraphics[scale=0.5]{dmrc-logo.eps}
\end{center}
\title{Slowcontrol for Y2L}
\author{Hyounggyu Kim} % Name of Authors
\affiliation{Center for Underground at Institute for Basic Science}
%\collaboration{}
%\date{\today}


\begin{abstract}
지하실험실 모니터링 전반에 대한 내용
\end{abstract}
%\keywords{KEYWORDS, KEYWORDS
%}
%\pacs{ PACS, PACS}
\maketitle
% main text

\tableofcontents

\section{모니터링 데이터 흐름}

\begin{figure}[htb]
\begin{center}
\includegraphics[width=0.9\textwidth]{diagrams/test.pdf}
\caption{This is a figure.}
\end{center}
\end{figure}

\section{컴퓨터 상태 정보 수집}

\subsection{Telegraf를 활용한 Windows 컴퓨터 모니터링}


다음과 같이 Telegraf를 윈도우즈 서비스로 등록할 수 있다.

\begin{verbatim}
.\telegraf.exe --service install \
               --config "\some\where\telegraf.conf"
\end{verbatim}

자세한 내용은 \href{https://www.influxdata.com/blog/using-telegraf-on-windows/}{이곳}을 참고.

\section{모니터링 프로그램 작성과 배포}



%\begin{table}
%\begin{center}
%\caption[TABLE\_CAPTION]{
%This is an example of table.
%}
%\begin{ruledtabular}
%\label{ref:tab1}
%\begin{tabular}{lccc}
%             &Sample1             & Sample2         &Sample3\\
%\hline
%BG1      & 780 $\pm$ 219& 311 $\pm$ 111 &19.5 $\pm$ 8.4\\
%\hline
%Background   &1038 $\pm$ 244& 420 $\pm$ 114 & 33.5 $\pm$ 11.4 \\
%Signal& 620 $\pm$ 83 &1541 $\pm$ 192 & 482 $\pm$ 74\\
%\hline
%Expected     &1658 $\pm$ 257&1961 $\pm$ 223 & 516 $\pm$ 75 \\
%Observed     &1712              & 1937              & 500 \\
%\end{tabular}
%\end{ruledtabular}
%\end{center}
%\end{table}

%Table~\ref{ref:tab1} is an example table. 
%\begin{cfigure1c}
%\begin{tabular}{cc}
%\includegraphics[width=0.43\textwidth]{dmrc-logo.eps} &
%\includegraphics[width=0.43\textwidth]{dmrc-logo.eps} \\
%(a) & (b) \\
%\end{tabular}
%\caption[Data fit]{
%An example of figure.
%}
%\label{ref:fig1}
%\end{cfigure1c}
%
%Figure~\ref{ref:fig1} is an example figure. 

\section{결론}

결론은 이러이러하다.

\section*{Acknowledgments}
Acknowledgments

%\begin{thebibliography}{99}

%\bibitem{bwlee} B.~W.~Lee and S.~Weiberg
%\Journal{\PRL}{39}{1977}{165}.

%\bibitem{Jungman} G.~Jungman, M.~Kamionkowski, K.~Griest,
%\Journal{\PREP}{267}{1996}{195}.

%\end{thebibliography}
\end{document}
